\section{Proving lemmas and propositions} \label{sec:lemmas_proofs}

\subsection{Stratified schemas}

Mixing unrestricted recursion with negation can make \emph{correct typing} 
 fail to find a unique semantics: 
 the same schema can admit several incompatible correct typings at once. 
 A standard fix, in the spirit of stratified negation in logic programming, 
 is to forbid negation from ever occurring on a recursive cycle 
 between shape labels. 
 This is captured by a dependency graph on shape labels with two kinds 
 of edges: \lean{DepPos S L1 L2} holds when the shape label L2 occurs 
 as a shape reference in L1's definition without being under a NOT (a positive dependency), 
 and \lean{DepNeg S L1 L2} holds when L2 occurs under some NOT (a negative dependency).

\begin{leanlisting}
abbrev DepPos (S : Schema) (L1 L2 : Label) : Prop := L2 ∈ (S.shapeDef L1).posDeps

abbrev DepNeg (S : Schema) (L1 L2 : Label) : Prop := L2 ∈ (S.shapeDef L1).negDeps
\end{leanlisting}

A schema with stratified negation admits a stratification \lean{Strat}, 
 represented as a structure that packages the number of strata \lean{k}, 
 the stratum \lean{str} of each label, 
 and proofs that this assignment respects the schema's dependency graph: 
 negative dependencies must strictly decrease the stratum (\lean{neg_strict}), 
 while positive dependencies may never increase it (\lean{pos_le}).

\begin{leanlisting}
structure Strat (S : Schema) where
  k : ℕ
  str : Label → ℕ
  mem_range : ∀ L ∈ S.labels, 1 ≤ str L ∧ str L ≤ k
  neg_strict : ∀ L1 ∈ S.labels, ∀ L2, DepNeg S L1 L2 →
                 str L2 < str L1
  pos_le : ∀ L1 ∈ S.labels, ∀ L2, DepPos S L1 L2 →
                 str L2 ≤ str L1
\end{leanlisting}

\lean{S} is a schema with stratified negation if it admits a stratification

\begin{leanlisting}
def Schema.StratNeg (S : Schema) : Prop := Nonempty (Strat S)
\end{leanlisting}

\subsection{Maximal Correct typing}

Lemma 1 covers the base case where negation is entirely absent: 
  for a negation-free schema \lean{S} (\lean{Schema.NegFree}, 
  every label's definition uses no NOT) 
  and any graph \lean{G}, 
  there is a correct typing that contains every other correct 
  typing of \lean{G} by \lean{S}. 
  The reason this case is easy is monotonicity: without negation, 
  \lean{Satisfies} only grows with the typing — 
  a node satisfying a shape expression under some typing keeps satisfying 
  it under any larger one — so simply taking the union of all correct typings (\lean{Schema.maximalTyping}) 
  is itself correct, 
  and contains every other one by construction. 

\begin{leanlisting}
theorem lemma1 (S : Schema) (hnf : S.NegFree) (G : Graph) :
    ∃ typingG : Typing, typingG ⊆ S.univTyping G ∧ S.CorrTyping G typingG ∧
      ∀ t', t' ⊆ S.univTyping G → S.CorrTyping G t' → t' ⊆ typingG
\end{leanlisting}

It is possible to construct a stratified-maximal correct typing \lean{MTyping}
  one stratum at a time. 

\begin{leanlisting}
def MTyping (G : Graph) (S : Schema) (str : Strat S) : Typing :=
  typ S str.str G str.k

def typ (S : Schema) (str : Label → ℕ) (G : Graph) : ℕ → Typing
  | 0 => ∅
  | i + 1 => maxAgree (S.stratRestrict str (i + 1)) G
                      (S.stratum str i) (typ S str G i)
\end{leanlisting}

where \lean{maxAgree S' G Lo fixed} is the union of all correct typings 
 of \lean{G} by \lean{S'} that agree with a given typing fixed on the labels in \lean{Lo} — 
 the per-stratum construction that \lean{typ}'s inductive step performs, 
 generalizing Lemma 1's maximal correct typing (the special case \lean{Lo = ∅}, 
 no agreement constraint).

For any schema \lean{S}, any stratification \lean{str} of \lean{S}
and any graph \lean{G}, the \lean{MTyping G S str}
is a correct typing of \lean{G} by \lean{S}.

\begin{leanlisting}
theorem proposition1 (S : Schema) (str : Strat S)
                    (hClosed : S.Closed) (G : Graph) :
    S.CorrTyping G (MTyping G S str)
\end{leanlisting}

The Lean proof strengthens the statement before inducting: 
 rather than reasoning about MTyping directly, 
 it shows by induction on the stratum index $i$ that 
 \lean{typ S str G i} satisfies two invariants at once — 
 it stays within $S_i$'s universe of typings, 
 and it is already a correct typing of $S_i$. 
 The base case ($i = 0$) is vacuous, since the initial stratum $L_0$ is empty. 
 The inductive step reduces to a single correctness fact about 
 \lean{maxAgree} (the union-of-typings construction underlying \lean{typ}), 
 applied with the previous stratum as the 
 "agreement set" — well-formedness of the schema together with the stratification's 
 \lean{neg_strict} condition is exactly what guarantees that $S_{i+1}$'s 
 negative dependencies stay confined there. 
 Proposition 1 is then just this invariant read off at the last stratum, $i = k$.

 \lean{hClosed} is a well-formedness hypothesis with no counterpart 
 in the paper's own statement of Proposition 1: 
 it requires every shape label referenced inside \lean{S}'s definitions 
 to itself be declared. 
 In the paper this holds automatically, 
 since schemas are only ever defined on their own label set; 
 in Lean, \lean{shapeDef} must be a total function, 
 so this closure condition has to be assumed explicitly.


 Lemma 2 states that for any two stratifications \lean{str1}, \lean{str2} 
  of a schema \lean{S} the maximal correct typing does not depend on 
  which stratification was chosen:

 \begin{leanlisting}
 theorem lemma2 (S : Schema) (hClosed : S.Closed) 
               (str1 str2 : Strat S)
               (G : Graph) : MTyping G S str1 = MTyping G S str2
 \end{leanlisting}

 Here the paper and the Lean development diverge more than for Proposition 1: 
 the paper's own proof is stated only as an idea, 
 asserting without further justification that a canonical \emph{most refined} 
 stratification exists which every other stratification agrees with. 
 The Lean proof has to supply this argument in full. 
 It introduces a canonical process that starts from the empty typing 
 and repeatedly reveals the least-stratum batch of labels still remaining; 
 an \emph{Independence} argument then shows, 
 by strong induction on the number of labels left, 
 that the process's final result does not depend on the order 
 in which batches are revealed — 
 the technical heart being a \emph{Swap} lemma letting 
 two mutually-independent batches be revealed in either order. 
 Running this process along the checkpoints of an arbitrary 
 stratification shows it agrees with \lean{MTyping} 
 for that stratification; 
 applying this to \lean{str1} and \lean{str2} 
 separately, both equal the same canonical value, giving Lemma 2.